\chapter{Linear Algebra Foundations}
\label{ch:linalg}

\section{Why this chapter matters}
This chapter is not a general linear algebra course. Its purpose is to develop, with enough precision to be useful later, the specific structures on which post-quantum cryptography is built: vector spaces over a finite field, linear maps and their matrices, and the rank--nullity relationship that governs when a linear system can be solved. The central object of lattice cryptography, the equation
\[
A\mathbf{s} + \mathbf{e} = \mathbf{b},
\]
is a linear system deliberately perturbed by a small error $\mathbf{e}$. To understand why the unperturbed system is easy and the perturbed one is (conjecturally) hard, we first need to be precise about what ``solving a linear system'' means.

\section{Learning goals}
After reading this chapter, you should be able to:

\begin{itemize}
  \item state the axioms of a vector space and recognise $\mathbb{F}_q^n$ as the working example;
  \item define linear independence, span, basis, and dimension, and know why dimension is well defined;
  \item describe a matrix as the coordinate form of a linear map, and define its rank, kernel, and image;
  \item state the rank--nullity theorem and use it to reason about solvability of $A\mathbf{x}=\mathbf{b}$;
  \item explain precisely why $A\mathbf{s}=\mathbf{b}$ is easy while $A\mathbf{s}+\mathbf{e}=\mathbf{b}$ is the hard LWE problem.
\end{itemize}

\section{Scalars: the underlying field}
Every vector space is built over a set of \emph{scalars} that can be added, subtracted, multiplied, and divided. That set is a \emph{field} (Chapter~\ref{ch:sagemath}, where the finite field $\mathbb{F}_q = \mathrm{GF}(q)$ was introduced). We recall the definition because the rest of the chapter depends on it.

\begin{definition}[Field]
A \emph{field} is a set $\mathbb{F}$ with two operations $+$ and $\cdot$ such that $(\mathbb{F},+)$ is an abelian group with identity $0$, the nonzero elements $(\mathbb{F}\setminus\{0\},\cdot)$ form an abelian group with identity $1$, and multiplication distributes over addition. In particular every nonzero element has a multiplicative inverse.
\end{definition}

The familiar fields are $\mathbb{Q}$, $\mathbb{R}$, and $\mathbb{C}$. Cryptography instead works over the finite field $\mathbb{F}_q$; for a prime $q$ this is simply the integers modulo $q$,
\[
\mathbb{F}_q = \mathbb{Z}/q\mathbb{Z} = \{0,1,\dots,q-1\},
\]
with arithmetic reduced modulo $q$. The property that makes it a field, rather than merely a ring, is that division is available: because $q$ is prime, every nonzero residue has an inverse. This is exactly what lets Gaussian elimination run, and it is why cryptographic parameters such as $q=3329$ (Kyber) are chosen prime.

\section{Vectors and vector spaces}
In school a vector is a coordinate tuple such as $(1,2)$. For our purposes a vector is an element of a \emph{vector space}: a structured container of field elements that can be added and scaled.

\begin{definition}[Vector space]
A \emph{vector space} over a field $\mathbb{F}$ is a set $V$ with an addition $V\times V\to V$ and a scalar multiplication $\mathbb{F}\times V\to V$ such that $(V,+)$ is an abelian group and, for all $a,b\in\mathbb{F}$ and $\mathbf{u},\mathbf{v}\in V$,
\[
a(\mathbf{u}+\mathbf{v}) = a\mathbf{u}+a\mathbf{v},\quad
(a+b)\mathbf{v} = a\mathbf{v}+b\mathbf{v},\quad
(ab)\mathbf{v} = a(b\mathbf{v}),\quad
1\mathbf{v} = \mathbf{v}.
\]
\end{definition}

\begin{example}[The space $\mathbb{F}_q^n$]
The set of column vectors
\[
\mathbb{F}_q^n = \left\{ (v_1,\dots,v_n)^\top : v_i \in \mathbb{F}_q \right\},
\]
with entrywise addition and scalar multiplication, is a vector space over $\mathbb{F}_q$. It is finite: it contains exactly $q^n$ elements. A Kyber secret key lives in $\mathbb{F}_{3329}^{256}$, so it is one of $3329^{256}$ possible vectors -- already far more than the number of atoms in the observable universe.
\end{example}

A typical secret is written
\[
\mathbf{s} = \begin{bmatrix}2\\5\\1\\7\end{bmatrix} \in \mathbb{F}_{q}^{4},
\qquad\text{or in practice}\qquad
\mathbf{s} \in \mathbb{F}_{3329}^{256}.
\]

\subsection{Sage experiment: creating a vector}

\begin{lstlisting}[language=Python]
F = GF(17)
v = vector(F, [1, 2, 3])
v
\end{lstlisting}

To inspect its type, use \texttt{type(v)}. Sage reports it as a \emph{free module element} rather than a plain list. This is not pedantry: a vector space over a field is the special case of a \emph{module} over a ring in which the ring happens to be a field, and Module-LWE (Chapter~\ref{ch:mlwe}) will use exactly the more general object.

\subsection{Vector addition and scalar multiplication}

\begin{lstlisting}[language=Python]
a = vector(F, [1, 2, 3])
b = vector(F, [5, 6, 7])
a + b        # (6, 8, 10)
3 * a        # (3, 6, 9)
\end{lstlisting}

All arithmetic is modulo the field characteristic: over $\mathbb{F}_7$ the sum $6+5$ is $4$, not $11$. Over $\mathbb{R}$ or $\mathbb{Q}$ addition has the familiar parallelogram picture of Figure~\ref{fig:vector-addition}.

\begin{figure}[H]
\centering
\begin{tikzpicture}[scale=0.9]
  \draw[pqcgray!40, thin] (-0.6,0) -- (4.8,0);
  \draw[pqcgray!40, thin] (0,-0.6) -- (0,3.6);
  \draw[pqcgray!40, dashed] (3,1) -- (4,3);
  \draw[pqcgray!40, dashed] (1,2) -- (4,3);
  \draw[basisvec] (0,0) -- (3,1);
  \draw[altvec] (0,0) -- (1,2);
  \draw[shortvec] (0,0) -- (4,3);
  \node[pqcblue, font=\small] at (3.35,0.7) {$u$};
  \node[pqcteal, font=\small] at (0.55,2.3) {$v$};
  \node[pqcorange, font=\small] at (4.35,2.85) {$u+v$};
\end{tikzpicture}
\caption{Vector addition by the parallelogram rule, for $u=(3,1)$ and $v=(1,2)$. Both vectors are drawn from the origin; completing the parallelogram (dashed sides) shows that the sum $u+v=(4,3)$ is exactly its diagonal. Coordinate-wise, this is the same rule the Sage experiment above computes entrywise: $(1,2,3)+(5,6,7)=(6,8,10)$.}
\label{fig:vector-addition}
\end{figure}

\section{Linear combinations, independence, and span}
The single most important operation in the whole subject is forming a \emph{linear combination}.

\begin{definition}[Linear combination and span]
Given vectors $\mathbf{v}_1,\dots,\mathbf{v}_k \in V$ and scalars $c_1,\dots,c_k\in\mathbb{F}$, the vector $\sum_{i=1}^k c_i\mathbf{v}_i$ is a \emph{linear combination} of the $\mathbf{v}_i$. The set of all such combinations is their \emph{span}, written $\mathrm{span}(\mathbf{v}_1,\dots,\mathbf{v}_k)$; it is the smallest subspace of $V$ containing every $\mathbf{v}_i$.
\end{definition}

\begin{definition}[Linear independence]
The vectors $\mathbf{v}_1,\dots,\mathbf{v}_k$ are \emph{linearly independent} if the only linear combination equal to the zero vector is the trivial one:
\[
\sum_{i=1}^k c_i \mathbf{v}_i = \mathbf{0} \quad\Longrightarrow\quad c_1 = \cdots = c_k = 0.
\]
Otherwise they are \emph{linearly dependent}, meaning one of them is a linear combination of the others.
\end{definition}

For instance $(1,0)$ and $(0,1)$ are independent, whereas $(1,2)$ and $(2,4)$ are dependent, since the second is twice the first.

\subsection{Sage experiment: detecting dependence through rank}
Dependence is detected numerically by \emph{rank}, defined precisely below: stacking the vectors as rows, the rank counts how many are genuinely independent.

\begin{lstlisting}[language=Python]
A = Matrix(GF(17), [[1, 2], [2, 4]])
A.rank()      # 1: only one independent direction
\end{lstlisting}

\section{Basis and dimension}
\begin{definition}[Basis]
A \emph{basis} of a vector space $V$ is a set of vectors that is both linearly independent and spans $V$. Equivalently, every $\mathbf{v}\in V$ has a \emph{unique} expression as a linear combination of the basis vectors.
\end{definition}

The standard basis of $\mathbb{F}_q^2$ is $\{(1,0),(0,1)\}$, since $(a,b)=a(1,0)+b(0,1)$. But $\{(2,1),(1,1)\}$ is an equally valid basis of the same space: the space is intrinsic, the basis is a chosen coordinate system for it. The number of vectors in a basis is not a matter of choice.

\begin{theorem}[Dimension is well defined]
Any two bases of a finite-dimensional vector space $V$ have the same number of elements. This common number is the \emph{dimension} $\dim V$.
\end{theorem}

Thus $\dim \mathbb{F}_q^n = n$. The freedom to change basis without changing the space is the seed of the entire theory of lattice reduction (Chapter~\ref{ch:lll}): a lattice, too, has many bases, and finding a \emph{good} one is the central algorithmic problem.

\subsection{Sage experiment: two bases for one space}

\begin{lstlisting}[language=Python]
B1 = Matrix(GF(17), [[1, 0], [0, 1]])
B2 = Matrix(GF(17), [[2, 1], [1, 1]])
# both have rank 2, so each is a basis of GF(17)^2
print(B1.rank(), B2.rank())
\end{lstlisting}

\section{Matrices as linear maps}
A matrix is best understood not as a grid of numbers but as the coordinate form of a \emph{linear map}.

\begin{definition}[Linear map]
A map $T\colon \mathbb{F}^n \to \mathbb{F}^m$ is \emph{linear} if it respects both operations:
\[
T(\mathbf{u}+\mathbf{v}) = T(\mathbf{u})+T(\mathbf{v}),\qquad T(a\mathbf{v}) = a\,T(\mathbf{v}).
\]
Once bases are fixed, every linear map is given by an $m\times n$ matrix $A$ via $T(\mathbf{v})=A\mathbf{v}$, and every such matrix defines a linear map.
\end{definition}

The product $A\mathbf{v}$ is itself a linear combination: it is the combination of the \emph{columns} of $A$ weighted by the entries of $\mathbf{v}$. Concretely, entry $i$ of the result is the inner product of row $i$ of $A$ with $\mathbf{v}$, which is exactly the double loop of Algorithm~\ref{alg:matvec}.

\begin{lstlisting}[language=Python]
F = GF(17)
A = Matrix(F, [[1, 2], [3, 4]])
v = vector(F, [5, 6])
A * v
\end{lstlisting}

\begin{algorithm}[H]
\caption{Matrix--Vector Product}
\label{alg:matvec}
\begin{algorithmic}[1]
\Require Matrix $A \in \mathbb{F}^{m\times n}$ with entries $A_{i,j}$; vector $v \in \mathbb{F}^n$
\Ensure $w = Av \in \mathbb{F}^m$
\For{$i = 1,\dots,m$}
    \State $w_i \gets 0$
    \For{$j = 1,\dots,n$}
        \State $w_i \gets w_i + A_{i,j}\, v_j$
    \EndFor
\EndFor
\State \Return $w$
\end{algorithmic}
\end{algorithm}

This same loop is foundational for the whole book: it computes $A\mathbf{s}$ in the LWE equation, and in ring form (Chapter~\ref{ch:rlwe}) it reappears as the schoolbook polynomial multiplication used by RLWE and Kyber.

\section{Rank, kernel, and image}
Three subspaces measure what a linear map $A\colon\mathbb{F}^n\to\mathbb{F}^m$ does.

\begin{definition}[Image, kernel, rank]
The \emph{image} (or column space) is $\operatorname{im}A = \{A\mathbf{v} : \mathbf{v}\in\mathbb{F}^n\}\subseteq\mathbb{F}^m$; its dimension is the \emph{rank} of $A$. The \emph{kernel} (or null space) is $\ker A = \{\mathbf{v}\in\mathbb{F}^n : A\mathbf{v}=\mathbf{0}\}\subseteq\mathbb{F}^n$; its dimension is the \emph{nullity}.
\end{definition}

The image is exactly the set of right-hand sides $\mathbf{b}$ for which $A\mathbf{x}=\mathbf{b}$ has a solution. The kernel measures how far a solution is from unique: if $A\mathbf{x}_0=\mathbf{b}$, then the full solution set is $\mathbf{x}_0 + \ker A$. These two dimensions are tied together.

\begin{theorem}[Rank--nullity]
For any linear map $A\colon \mathbb{F}^n \to \mathbb{F}^m$,
\[
\dim(\ker A) + \operatorname{rank}(A) = n.
\]
\end{theorem}

In particular, a square matrix $A\in\mathbb{F}^{n\times n}$ is invertible if and only if its rank is $n$, equivalently $\ker A=\{\mathbf{0}\}$, equivalently its columns are a basis of $\mathbb{F}^n$. When $A$ is invertible, $A\mathbf{x}=\mathbf{b}$ has exactly one solution for every $\mathbf{b}$.

\subsection{Sage experiment: rank governs invertibility}

\begin{lstlisting}[language=Python]
F = GF(17)
for _ in range(10):
    A = random_matrix(F, 4, 4)
    print(A.rank())
\end{lstlisting}

For a uniformly random $n\times n$ matrix over $\mathbb{F}_q$, the invertibility probability is
\[
\prod_{i=1}^{n}(1-q^{-i}).
\]
Thus the claim depends on $q$: for this $4\times4$ example over $\mathbb{F}_{17}$ it is high, while over $\mathbb{F}_2$ it is only about $30.8\%$.

\section{Solving linear systems, and the leap to LWE}
Given $A$ and $\mathbf{b}$, the equation $A\mathbf{x}=\mathbf{b}$ is solved by \emph{Gaussian elimination}: row operations reduce $A$ to echelon form, and back-substitution recovers $\mathbf{x}$.

\begin{proposition}[Linear systems are easy]
Over a field, Gaussian elimination solves an $n\times n$ system $A\mathbf{x}=\mathbf{b}$ (or reports that no solution exists) using $O(n^3)$ field operations. If $A$ is invertible the solution is unique.
\end{proposition}

This efficiency is precisely why a plain linear system carries \emph{no} cryptographic value: anyone who sees $A$ and $\mathbf{b}=A\mathbf{s}$ can recover the secret $\mathbf{s}$ in cubic time. LWE changes the setting. Given a public, typically rectangular random matrix $A\in\mathbb{Z}_q^{m\times n}$ with $m>n$, a secret $\mathbf{s}$ drawn from a prescribed small distribution, and an independent small error vector $\mathbf{e}$, one observes samples
\[
\mathbf{b} = A\mathbf{s} + \mathbf{e},
\]
with arithmetic modulo $q$. Unlike an invertible square system, this overdetermined system will generally have no exact solution for a nonzero random error. Gaussian elimination can detect inconsistency, but it does not recover the planted secret from the noisy samples. Recovering $\mathbf{s}$ is the \emph{Learning With Errors} problem (Chapter~\ref{ch:lwe}).

\begin{remark}
It is worth pausing on how little changes on the page and how much changes in difficulty. For an invertible square matrix, $\mathbf{s}\mapsto A\mathbf{s}$ is easy to invert. With many noisy samples, the observed value also depends on independently sampled error; cryptographic hardness lives in distinguishing and exploiting that structure.
\end{remark}

\section{Experiments}

\begin{lstlisting}[language=Python]
F = GF(17)
# 1. random vectors
V = VectorSpace(F, 4)
for _ in range(5): print(V.random_element())

# 2. a matrix acting on a vector
A = random_matrix(F, 6, 4); v = random_vector(F, 4)
print(A * v)

# 3. exact solve is easy; try recovering s from b = A*s
s = random_vector(F, 4); b = A * s
print(A.solve_right(b) == s)   # True when A has full column rank

# 4. perturb an overdetermined system: it normally becomes inconsistent
e = vector(F, [1, 0, 0, 16, 0, 0])   # 16 == -1 mod 17
try:
    A.solve_right(b + e)
except ValueError:
    print("no exact solution: noisy samples are inconsistent")
\end{lstlisting}

The experiment is a toy illustration only: real LWE specifies the dimensions and the secret and error distributions precisely.

\section{Connection to the next chapter}
The next chapter, \emph{Polynomial Rings}, keeps the linear-algebraic viewpoint but changes the scalars-and-vectors picture into a single richer object. A structured (circulant) matrix will be encoded by one polynomial, matrix--vector multiplication will become polynomial multiplication, and the vector space $\mathbb{F}_q^n$ will be upgraded to the quotient ring
\[
R_q = \mathbb{Z}_q[x]/(x^n + 1),
\]
the natural home of RLWE and Kyber.
